\section*{Epistemological Manifesto: Paradigm Shift and Natural Philosophy}
\addcontentsline{toc}{section}{Epistemological Manifesto: Paradigm Shift and Natural Philosophy}
\label{sec:epistemolojik_manifest}

\subsection*{1. Epistemological Parsimony and the Conformal Geometric Perspective}

The foundational questions encountered in modern theoretical physics and cosmology—such as the nature of dark matter, dark energy, cosmic inflation, and the interface between gravity and quantum scales—motivate the exploration of alternative geometric formulations:
\begin{itemize}
    \item Galactic rotation curves diverging from Keplerian dynamics led to the postulation of cold dark matter.
    \item Distance-redshift relations of Type Ia supernovae led to the hypothesis of dark energy and cosmic acceleration.
    \item Horizon and flatness puzzles motivated cosmic inflation models in the early universe.
    \item Mass generation and confinement problems in field theory continue to stimulate active non-perturbative investigation.
\end{itemize}

This treatise investigates a conformal framework wherein, rather than metric spatial expansion, local atomic and nuclear length standards undergo secular conformal contraction ($r_{\mathrm{atom}}(t) \to 0$) against a static Euclidean background ($\mathcal{E}_\infty$). In the spirit of epistemological parsimony and Occam's razor, our objective is to examine how far scale-covariant differential geometry and conformal duality can account for observed phenomena with minimal axiomatic overhead, offering an exploratory geometric perspective on open cosmological problems.

\subsection*{2. Mathematical Legitimacy: Conformal Duality and Nilpotent Algebra}

The equivalence between spatial expansion and atomic ruler contraction is not merely an interpretative preference; it is a rigorous theorem of differential geometry. In General Relativity, the conformal transformation of the metric tensor:
\begin{equation}
    g_{\mu\nu} \to \Omega^2(t) g_{\mu\nu}
\end{equation}
establishes an exact geometric isomorphism (\textit{conformal duality}) between coordinate systems. As noted by theorists such as Christof Wetterich (2013), Fred Hoyle, Jayant Narlikar, and Roger Penrose across various contexts, an expanding metric space ($a(t) \propto e^{H_0 t}$) and a static background with contracting atomic rulers ($\lambda(t) \propto e^{-H_0 t}$) generate mathematically identical physical observables.

The dual nilpotent algebra ($\kappa^2 = 0$) and the Scale-Covariant Derivative $\mathbb{D}_t \equiv \frac{d}{dt} + \kappa H_0$ developed in this work elevate this conformal symmetry from a passive coordinate gauge choice into a causal and dynamical physical ontology. While space retains its quiescent Euclidean structure, matter is defined as a dynamic Scalor contracting along its internal phase trajectory.

\subsection*{3. Physical Foundations of Fundamental Constants: $c$ and the Sommerfeld Constant $\alpha$}

Standard physics accepts fundamental dimensional and dimensionless constants as empirical givens fixed at the Big Bang without underlying physical mechanisms:
\begin{itemize}
    \item \textit{``Why is the speed of light approximately $3 \times 10^8\text{ m/s}$, and why does it serve as an impassable velocity ceiling?''}
    \item \textit{``What is the physical meaning of the Sommerfeld fine-structure constant $\alpha$ in atomic kinematics?''}
\end{itemize}

Within our hydro-geometric framework:
\begin{enumerate}
    \item \textbf{The Speed of Light ($c$):} Far from being an arbitrary velocity limit, $c$ is the transverse elastic phase transmission speed of the quiescent vacuum continuum ($v_{\mathrm{phase}} = 1/\sqrt{\varepsilon_0\mu_0} = c$) and the characteristic radial contraction phase ceiling of matter vortices ($v_c = c$).
    \item \textbf{The Fine-Structure Constant ($\alpha \approx 1/137.036$):} Represents the \textbf{geometric vortex pitch ratio} required for an orbital electron to balance centrifugal acceleration against the nucleus's radial contraction phase ceiling ($v_r = v_c = c$) without radiative collapse:
          \begin{equation}
              \alpha \equiv \frac{v_\theta}{v_r} = \frac{v_{\mathrm{Bohr}}}{c} = \frac{e^2}{4\pi\varepsilon_0\hbar c} \approx \frac{1}{137.036}.
          \end{equation}
          The relativistic Dirac threshold ($Z\alpha \to 1$) marks the boundary where tangential orbital velocity reaches the contraction ceiling ($v_\theta = c$), inducing Schwinger vacuum polarization ($e^+ e^-$ production) to stabilize the effective nuclear charge.
\end{enumerate}
Reducing these fundamental parameters to kinematic ratios resolves an enduring ontological puzzle in theoretical physics.

\subsection*{4. Invariance of Dimensionless Ratios and Local Laboratory Tests (Eötvös / MICROSCOPE)}

The compatibility of the Unified Contraction Hypothesis with local high-precision laboratory experiments is guaranteed by the scale-invariance of dimensionless ratios. If atomic length scales contracted independently of mass scaling, satellite experiments like MICROSCOPE and terrestrial Eötvös balances would detect differential accelerations, falsifying the theory.

Because all bound mass states acquire conformal mass synchronously according to $m(t) = m_0 e^{H_0 t}$, the measured ratio of any two laboratory masses remains strictly constant:
\begin{equation}
    \frac{d}{dt} \left( \frac{m_{\mathrm{test}}}{m_{\mathrm{ref}}} \right) = \frac{d}{dt} \left( \frac{m_{1,0} e^{H_0 t}}{m_{2,0} e^{H_0 t}} \right) \equiv 0.
\end{equation}
Similarly, the fine-structure constant $\alpha$ and the proton-to-electron mass ratio $\mu = m_p/m_e$ are time-invariant ($\dot{\alpha}/\alpha = 0, \dot{\mu}/\mu = 0$). This explains why tests of the Weak Equivalence Principle (WEP) confirm the absence of differential acceleration down to $\eta < 10^{-15}$.

\subsection*{5. Unifying Celestial and Terrestrial Dynamics: Virial Decoupling and Ephemerides}

Cosmological models often restrict their scope to intergalactic scales, detached from local Solar System constraints. The Unified Contraction Hypothesis seamlessly connects deep-space phenomenology with local celestial mechanics:
\begin{itemize}
    \item \textbf{Local Virial Decoupling and LLR Measurements:} In gravitationally bound systems dominated by local potential ($\nabla^2\Phi \gg 3H_0^2$), orbits decouple from cosmological background expansion in accordance with the McVittie theorem. The measured $3.82\text{ cm/yr}$ recession of the Moon via Lunar Laser Ranging (LLR) is driven by ocean tidal friction, as confirmed by satellite altimetry (TOPEX/Poseidon).
    \item \textbf{Ephemeris and Radar Telemetry Limits:} Because ruler contraction modulates atomic transition frequencies rather than spatial coordinates, the theory remains strictly compatible with radar telemetry constraints on planetary ephemerides ($< 1.5\text{ cm/yr}$ per AU).
\end{itemize}

\subsection*{6. Conceptual Clarity: ``Vacuum Fluid'' vs. 19th-Century Aether (Emergent Spacetime)}

The terminology employed herein---\textit{vacuum flow}, \textit{vortex}, and \textit{Bernoulli gradient}---must not be conflated with 19th-century mechanical luminiferous aether hypotheses. Classical aether was conceived as a dissipative, particulate substance, a concept decisively refuted by the Michelson--Morley experiment.

Here, the background continuum $\mathcal{E}_\infty$ is a frictionless, non-dissipative, massless, and smooth Euclidean manifold. The vacuum flow induced by contracting matter aligns directly with modern formulations of \textit{Emergent Spacetime} and \textit{Analogue/Acoustic Gravity}. Particles are not rigid spheres suspended in space; they are localized topological vortex solitons within the field continuum.

\subsection*{7. Computational Verification and Open-Source Solvers: Minimal CMB Pipeline}

The analytical framework is accompanied by an open-source, reproducible Python solver package (\texttt{scale-cmb-solver}, available at \href{https://github.com/FurkanHaydari/universe-is-not-expanding}{\texttt{FurkanHaydari/universe-is-not-expanding}}) that models the full Planck 2018 CMB angular power spectrum ($\ell = 2\text{--}2500$):
\begin{enumerate}
    \item \textbf{Planck CMB Angular Power Spectrum Solver:} The solver (\texttt{analysis/cmb\_spectrum\_analysis.py}) incorporates the acoustic standing wave transfer equation of contracting plasma, baryon drag ($R_b \approx 0.058$), the acoustic peak spacing ($\Delta\ell \approx 293.5$), early ISW enhancement, and Silk diffusion damping ($\ell_D \approx 1638.7$). It fits the Planck 2018 temperature anisotropy ($D_\ell^{\mathrm{TT}}$) binned data points with high fidelity, reproducing the first three acoustic peaks at $\ell_1 = 220.9$ (observed: 220.0), $\ell_2 = 518.4$ (observed: 537.5), and $\ell_3 = 820.2$ (observed: 810.8) along with their relative amplitude ratios ($H_1/H_2 = 2.169$, $H_3/H_2 = 0.886$).
    \item \textbf{Deuterium and Primordial Abundances Simulation:} The observed intergalactic deuterium abundance ($D/H \sim 2.5 \times 10^{-5}$) is verified through a multi-component \textit{Astrophysical Nuclear Network Simulator} balancing stellar astration with photodisintegration and spallation yields in the quiescent medium.
\end{enumerate}

\subsection*{8. Sociology of Science and Institutional Resistance}

As Max Planck observed: \textit{``A new scientific truth does not triumph by convincing its opponents and making them see the light, but rather because its opponents eventually die, and a new generation grows up that is familiar with it.''}

The primary obstacle facing this framework is the institutional inertia of established paradigms:
\begin{itemize}
    \item Eliminating dark matter and dark energy challenges extensive research programs and funding allocated over decades.
    \item Replacing the Big Bang ($t=0$) singularity with an asymptotic contraction phase transition necessitates a reevaluation of literature founded on expanding metrics.
\end{itemize}
Accordingly, polemical discourse is set aside in favor of an axiomatic, mathematically complete, and empirically testable physical architecture.

\vspace{1.5em}
\hrule
\vspace{1.5em}
